% !PRINTED 24
\weberpage{24}
%<<CONT
\begingroup
\setlength{\abovedisplayskip}{5pt plus 2pt minus 1pt}
\setlength{\belowdisplayskip}{5pt plus 2pt minus 1pt}
\setlength{\abovedisplayshortskip}{3pt plus 1pt}
\setlength{\belowdisplayshortskip}{4pt plus 1pt}
\linespread{0.94}\selectfont
equal to the highest of the degrees of the summands, and can decrease
only in the special case in which the highest degree occurs in several
summands and the sum of the coefficients of the highest powers is
zero.

In the multiplication of two or more whole rational functions there
arises a function whose degree is equal to the sum of the degrees of
the factors.

To survey the law of formation of the coefficients for the product,
put
\begin{equation}
\begin{aligned}
 A(x)&=a_0x^m+a_1x^{m-1}+a_2x^{m-2}+\cdots,\\
 B(x)&=b_0x^n+b_1x^{n-1}+b_2x^{n-2}+\cdots,
\end{aligned}
\tag{2}
\end{equation}
\begin{equation}
 A(x)B(x)=C(x)=c_0x^{m+n}+c_1x^{m+n-1}+c_2x^{m+n-2}+\cdots
\tag{3}
\end{equation}
and obtain
\[
\begin{aligned}
 c_0&=a_0b_0,\\
 c_1&=a_0b_1+a_1b_0,\\
 c_2&=a_0b_2+a_1b_1+a_2b_0,\\
 &\cdots,
\end{aligned}
\]
or, in general notation, if $\nu$ denotes one of the numbers
$0,1,2,\ldots,m+n$:
\begin{equation}
 c_\nu=a_0b_\nu+a_1b_{\nu-1}+a_2b_{\nu-2}+\cdots
       +a_{\nu-1}b_1+a_\nu b_0,
\tag{4}
\end{equation}
where all $a$ whose index is greater than $m$, and all $b$ whose index
is greater than $n$, are to be set equal to zero. For $c_\nu$ is the
coefficient of $x^{m+n-\nu}$, and hence $c_\nu x^{m+n-\nu}$ arises by
multiplying all terms of the form $a_\mu x^{m-\mu}$ by all terms of the
form $b_{\nu-\mu}x^{n-\nu+\mu}$ and then adding. If $a_0$ and
$b_0=1$---that is, if both are equal to one---then also $c_0=1$.%
% The source's compressed “a_0 and b_0=1” is retained symbol-for-symbol,
% while the translation makes its mathematically intended joint scope explicit.

\gesperrt{If in all factors of such a product the highest power of
$x$ has coefficient $1$, then in the product too the coefficient of
the highest power is $=1$.}

From these rules for calculating with whole functions it follows that
the rules of calculation expressed in formulas such as $ab=ba$,
$(ab)c=a(bc)$, $(a+b)c=ac+bc$ and the like also hold when $a,b,c$ are
whole functions of $x$, and in the sense that whole functions are to
be regarded as equal only when equal powers have equal coefficients.

By whole functions of several variables we understand sums of products
of whole, positive powers of the variables $x,y,z\ldots$ with arbitrary
coefficients that count as constant. One can think of these functions
as arising
\endgroup
%CONT>>
